\documentclass[a4paper,12pt]{ltjsarticle}
\usepackage[top=25mm,bottom=25mm,left=25mm,right=25mm]{geometry}
\usepackage{graphicx}
\usepackage{float}
\usepackage{booktabs}
\usepackage{array}
\usepackage{enumitem}
\usepackage{hyperref}
\usepackage{caption}

\begin{document}

\begin{center}
  {\Large \textbf{第10回議事録}}\\[8pt]
  {\large チーム4}
\end{center}

\vspace{4mm}

%-------------------------------------------------------------------
\section*{1.基本情報}
\begin{tabular}{ll}
  \toprule
  日時     & 2026年7月1日（水） \\
  会議場所 & 731教室 \\
  目的     & 進捗報告 \\
  チーム番号 & 4 \\
  議事録作成者 & 細澤悠真 \\
  \bottomrule
\end{tabular}

\vspace{4mm}
\noindent\textbf{参加者}

\begin{tabular}{ll}
  2年 & 山口汐里，佐藤夏美，成毛海人 \\
  3年 & 細澤悠真，武本龍，酒井睦基 \\
\end{tabular}

%-------------------------------------------------------------------
\section*{2.議題一覧}
\begin{enumerate}[label=議題\arabic*：]
  \item 議論して決まった内容
  \item 計画からの変更点
  \item 個人毎の次回までに実施する内容
\end{enumerate}

%-------------------------------------------------------------------
\section*{3.議題ごとの内容}

\subsection*{議題①：議論して決まった内容}
\subsubsection*{決定事項}
\begin{itemize}[leftmargin=1.5em]
  \item 輪唱システムと可動システムの結合が完了した。Master Arduino・Slave Arduino・プログラム部分のいずれも結合済みである。
  \item 可動部分の1周あたりの計測時間は約2200ms（10周で21.91秒）であることを測定により確認した。
  \item 結合後の動作仕様を以下のとおり決定した。
  \begin{itemize}
    \item ロードセルの値が閾値を超えたらフラグを立て，2200ms後に輪唱を開始する。
    \item ロードセルの値が閾値を下回ったらフラグを下げ，輪唱を終了する。
    \item BPMに依存して可動部分の回転速度を変化させる。
  \end{itemize}
  \item 結合テストを実施した結果，輪唱は実行可能であるが，可動部分の実行で不具合が発生した。不具合特定のための取り組みにより，ロードセルの異常値検知に対してデジタルピンを逆転させ，閾値を正常値へ修正した後，サーボモータの回転を確認した。
  \item I2C通信の転送レートをFast-mode（400\,kbps）へ変更し，1ビットあたりの送信時間を短縮して全体の転送遅延を削減した。
  \item BPM変化時の平滑化ロジック（移動平均等）は実装しないこととした（詳細は議題②を参照）。
\end{itemize}

\subsubsection*{議論内容}
\begin{itemize}[leftmargin=1.5em]
  \item 可動部分の不具合特定のため，以下を実施した。
  \begin{itemize}
    \item 回路のチェック。
    \item 電源の再起動。
    \item Arduinoを複数接続しているため，別のArduinoに別プログラムを書き込んでいないかの確認。
    \item ロードセルの値が異常値を検知していたため，デジタルピンを逆転させて閾値を正常値に修正し，サーボモータの回転を確認した。
  \end{itemize}
  \item 楽器ごとの接続不具合が残っており，個別の対応が必要である。
  \item PPM補正用ポテンショメータの視覚化について議論した。ポテンショメータをどれくらい回しているかをProcessing上に描画するなどして可視化する案が出た。
  \item 楽器ごとの音量調節が必要である。
  \item 補正方式について，小節ごとに補正をかけるためPLL補正（Phase Locked Loop：基準信号と出力信号の位相を比較しハードウェア的に発振周波数を同期・補正する手法）のほうがPPM補正（Parts Per Million：発振回路が持つ数ppm〜数十ppmの周波数誤差をソフトウェア的・数学的に補正する手法）よりも適しているとの見解が示された。
  \item BPM変化時の平滑化ロジックは未実装であり，移動平均などの平滑化ロジックを入れる可能性を検討した。ただし，浮動小数点数を扱うため処理速度低下による遅延が発生する懸念がある。
  \item サーボモータが電力不足により立ち上がりが遅く，別途電力供給が必要である。供給方法として，山口より「9Vアルカリ乾電池を別途供給電源にする」案，酒井より「電池ボックスから直接サーボモータに接続して電源供給する」案が出た。
\end{itemize}

\subsubsection*{ToDo / アクション}
\begin{itemize}[leftmargin=1.5em]
  \item 成果報告会の題目申請。
  \item 成果報告会のスライド作成。
  \item リラクゼーション効果・楽曲識別・楽器識別のためのアンケート作成（3年生が主体）。
  \item 【3年】可動部分の装飾，装飾の再調達，進捗状況の整理・マネジメント，結合テストの準備，GitHub ProjectのIssue更新。
  \item 【2年】7月15日の成果報告会に向けたスライド作成の開始，輪唱部分のユニットテスト，輪唱部分と可動部分の結合，輪唱プログラムと可動プログラムの結合，結合テスト，成果報告会の題目申請。
\end{itemize}

\subsubsection*{保留・次回検討事項}
\begin{itemize}[leftmargin=1.5em]
  \item サーボモータの供給電源方式（9Vアルカリ乾電池／電池ボックス直結）の選定。
  \item PPM補正用ポテンショメータの視覚化（Processing描画）の実装方法。
  \item 楽器ごとの音量調節および楽器ごとの接続不具合への対応。
\end{itemize}

%-------------------------------------------------------------------
\subsection*{議題②：計画からの変更点}

\subsubsection*{決定事項}
\begin{itemize}[leftmargin=1.5em]
  \item BPM変化時の平滑化補正アルゴリズムは実装しないこととした。測定の結果，そもそも補正が必要なほどの遅延が存在していなかったためである。
  \item 本件は測定によって判明した事象であるため，計画書の修正は行わない。
  \item ただし，チームの作業ログなどの変更は必要であり，TPMの修正を行う。
\end{itemize}

\subsubsection*{議論内容}
\begin{itemize}[leftmargin=1.5em]
  \item 酒井より，補正が必要なほどの遅延が存在していなかったため平滑化アルゴリズムの実装を行わない旨の共有があった。計画書自体の変更はないが，作業ログおよびTPMの修正が必要であるとの結論に至った。
\end{itemize}

\subsubsection*{ToDo / アクション}
\begin{itemize}[leftmargin=1.5em]
  \item TPMの修正……酒井
  \item チーム作業ログの更新。
\end{itemize}

\subsubsection*{保留・次回検討事項}
特になし。

%-------------------------------------------------------------------
\subsection*{議題③：個人毎の次回までに実施する内容}

\subsubsection*{決定事項}
各学年が次回までに実施する内容を以下に示す。
\begin{itemize}[leftmargin=1.5em]
  \item 【3年】リラクゼーション効果・楽曲識別・楽器識別のためのアンケート作成（3年生が主体となって進める），可動部分の装飾，装飾の再調達，進捗状況の整理・マネジメント，結合テストの準備，GitHub ProjectのIssue更新。
  \item 【2年】7月15日の成果報告会に向けたスライド作成の開始，輪唱部分のユニットテスト，輪唱部分と可動部分の結合，輪唱プログラムと可動プログラムの結合，結合テスト，成果発表会の題目申請。
\end{itemize}

\subsubsection*{議論内容}
\begin{itemize}[leftmargin=1.5em]
  \item 来週の成果報告に向けて，リラクゼーション効果および楽曲・楽器識別を評価するためのアンケートを作成する必要がある。当該作業は3年生が主体となって進める。
  \item 2年生には，7月15日の成果報告会に向けたスライド作成を開始してもらう。
  \item 次回までの最優先事項は，輪唱システムと可動システムの結合テストの完遂，および可動部分の不具合解消・装飾である。
\end{itemize}

\subsubsection*{ToDo / アクション}
議題①のとおり。

\subsubsection*{保留・次回検討事項}
特になし。

%-------------------------------------------------------------------
\section*{4.次回予定}
\begin{tabular}{ll}
  \toprule
  日時・場所 & 2026年7月8日（水）・731教室 \\
  議題       & システム評価の会（結合テスト結果報告，成果報告会準備）\\
  その他，特記事項 & 成果報告会は7月15日実施。アンケート作成は3年生主体，スライド作成は2年生が開始 \\
  \bottomrule
\end{tabular}

\end{document}
